\documentclass[lualatex,ja=standard]{bxjsarticle}
\usepackage{amsmath,amssymb}
\usepackage{enumitem}
\setlist[enumerate]{itemsep=0.6em, topsep=0.5em}
\pagestyle{empty}

\begin{document}

\noindent\textbf{数IA 共通テスト演習（図形・確率・データ分析） 解答}
\hrule
\vspace{0.8em}

\begin{enumerate}
\item
\begin{enumerate}[label=(\arabic*)]
\item 余弦定理より
\[
BC^2=AB^2+AC^2-2\cdot AB\cdot AC\cos A
\]
\[
=6^2+8^2-2\cdot6\cdot8\cos120^\circ
=36+64-96\left(-\frac12\right)=148
\]
\[
BC=2\sqrt{37}
\]
\item 面積公式より
\[
S=\frac12\cdot AB\cdot AC\sin A
=\frac12\cdot6\cdot8\sin120^\circ
=24\cdot\frac{\sqrt3}{2}
=12\sqrt3
\]
\end{enumerate}

\item 余弦定理より
\[
AC^2=AB^2+BC^2-2\cdot AB\cdot BC\cos B
\]
\[
8^2=5^2+7^2-2\cdot5\cdot7\cos B
\]
\[
64=74-70\cos B
\Rightarrow \cos B=\frac{1}{7}
\]
\[
\cos B>0
\]
より、$\angle B$ は鋭角。

\item 全事象は
\[
\binom92=36
\]
和が6の倍数になる組は
\[
(1,5),(1,11\text{不可}),(2,4),(2,10\text{不可}),(3,9),(4,8),(5,7)
\]
を1〜9内で整理すると
\[
(1,5),(2,4),(3,9),(4,8),(5,7)
\]
の5通り。
よって確率は
\[
\frac{5}{36}
\]

\item 全事象は
\[
\binom93=84
\]
「ちょうど2色」は補集合を用いる。
\[
\text{(1色のみ)}+\text{(3色すべて)}
\]
を全体から引く。

1色のみ:
\[
\binom43+\binom33+\binom23=4+1+0=5
\]

3色すべて:
\[
\binom41\binom31\binom21=24
\]

したがって
\[
84-(5+24)=55
\]
よって確率は
\[
\frac{55}{84}
\]

\item 条件「和が9以上」の場合の標本空間は
\[
(3,6),(4,5),(5,4),(6,3),
\]
\[
(4,6),(5,5),(6,4),
\]
\[
(5,6),(6,5),
\]
\[
(6,6)
\]
の10通り。

このうち積が奇数なのは両方奇数のときのみで
\[
(5,5)
\]
の1通り。
したがって積が偶数は9通り。

よって条件付き確率は
\[
\frac{9}{10}
\]

\item
\begin{itemize}
\item[(ア)] 中央値66点より、20番目と21番目の平均が66。必ずしも66点以上が20人以上とは限らない。
\item[(イ)] 必ずしも成り立たない。
\item[(ウ)] 平均64点より、最高点は64点以上であることが必要。
\end{itemize}
よって必ず正しいのは
\[
\text{(ウ)}
\]

\item データ（昇順）
\[
4,\ 5,\ 6,\ 7,\ 7,\ 8,\ 9,\ 10,\ 12,\ 12
\]

中央値（$Q_2$）は5番目と6番目の平均：
\[
Q_2=\frac{7+8}{2}=7.5
\]

下位5個
\[
4,\ 5,\ 6,\ 7,\ 7
\]
の中央値より
\[
Q_1=6
\]

上位5個
\[
8,\ 9,\ 10,\ 12,\ 12
\]
の中央値より
\[
Q_3=10
\]

四分位範囲は
\[
Q_3-Q_1=10-6=4
\]

\item 偏差値の公式
\[
\text{偏差値}=50+10\frac{x-\mu}{\sigma}
\]
を用いる（$\mu=58,\ \sigma=12$）。

得点74点の偏差値:
\[
50+10\cdot\frac{74-58}{12}
=50+\frac{40}{3}
=\frac{190}{3}\approx63.3
\]

偏差値65に相当する得点を $x$ とすると
\[
65=50+10\frac{x-58}{12}
\Rightarrow 15=10\frac{x-58}{12}
\Rightarrow x-58=18
\Rightarrow x=76
\]
\end{enumerate}

\end{document}
